\subsection*{Selected reference and incidence terminology}
\begin{description}[style=nextline]
\item[Executing reference.] A selected presentation at a declared registration interface, with retained local context and any orientation data required by the intended readout. It is part of the presentation specification, not a physically privileged frame.
\item[Situated account.] A record and readout indexed to an executing reference. The word subjective in the motivating discussion has this operational meaning, not the meaning inaccurate or unsigned.
\item[Reference signing.] The coordinates $\sigma_{t_*}$ on a two-member orientation torsor after one member $t_*$ has been selected. The normalization $\sigma_{t_*}(t_*)=+1$ is not a complexity, energy, or event-depth normalization.
\item[Incidence presentation.] An object together with the relation in which it is presented, such as $(u,B)$ with $u\in B$. Required phase or endpoint data must be included rather than inferred from the object alone.
\item[Orientation reconciliation.] A covariant comparison between the orientation torsors at an admitted incidence. In local coordinates its coefficient $\kappa$ obeys (IR5); its existence is not guaranteed by selecting references.
\item[Reference-independent statement.] A statement unchanged by the declared simultaneous transformation of its represented objects, reference data and comparison laws. This is a property on a specified domain, not a claim of an unrestricted physical frame symmetry.
\item[Phase-transport automorphism.] The pair $(\operatorname{Ad}_r,x\mapsto rx)$ on the registered-history action groupoid. Its action on a separately named history-state readout requires its own intertwining equation.
\item[Signed-axis extension.] The split operator group $\widetilde A\cong A_5\times C_2$, distinguished from the analyzer-sector group $S_5$, the common-cover graph $\mathsf K_{120}$, and the 120 bare face incidences.
\end{description}
